Esta sección presenta el desarrollo del mecanismo de almacenamiento comprimido basado en
Blosc incorporado en TMFQSfullstate. La implementación se apoya en una interfaz común de
almacenamiento del vector de estado, sobre la cual coexisten una realización no comprimida y
realizaciones comprimidas. La descripción se organiza alrededor de la disposición en bloques
del vector de estado, la caché LRU de bloques descomprimidos y la estrategia de acceso local
utilizada durante la simulación.

\begin{figure}[H]
    \centering
    \fbox{\parbox{0.82\linewidth}{\centering
        Espacio reservado para un diagrama de módulos del desarrollo realizado, mostrando la relación entre el registro cuántico, el esquema basado en Blosc,
        la disposición en bloques, la caché y el flujo de compresión/descompresión.}}
    \caption{Arquitectura del desarrollo implementado para el esquema basado en Blosc}
    \label{fig:development-architecture}
\end{figure}
